% !TeX program = lualatex
% Compile twice: lualatex 118.tex
% All references and prose are embedded; no BibTeX database or external inputs.
\documentclass[11pt,a4paper]{article}
\usepackage[margin=24mm,headheight=14pt]{geometry}
\usepackage{fontspec}
\setmainfont{Latin Modern Roman}
\setsansfont{Latin Modern Sans}
\setmonofont{Latin Modern Mono}
\usepackage{amsmath,amssymb,mathtools}
\usepackage{unicode-math}
\setmathfont{Latin Modern Math}
\usepackage{microtype}
\usepackage{enumitem,xurl,needspace}
\usepackage[dvipsnames]{xcolor}
\usepackage{hyperref}
\hypersetup{unicode=true,colorlinks=true,linkcolor=MidnightBlue,urlcolor=MidnightBlue,
 pdftitle={Research attempt: Problem 118},pdfauthor={Research notebook},bookmarksnumbered=true}
\usepackage{bookmark}
\usepackage{tocloft}
\setlength{\cftsecnumwidth}{3em}
\cftsetpnumwidth{2.5em}
\cftsetrmarg{4em}
\usepackage{titlesec}
\titleformat{\section}{\Large\bfseries\raggedright}{\thesection}{0.7em}{}
\usepackage{fancyhdr}
\pagestyle{fancy}\fancyhf{}
\fancyhead[L]{\small\sffamily Open Problems Math | Research notebook}
\fancyhead[R]{\small Problem 118 | 27 September 2026}
\fancyfoot[C]{\thepage}
\setlength{\parindent}{0pt}
\setlength{\parskip}{3pt plus 1pt}
\setlength{\emergencystretch}{3em}
\setcounter{tocdepth}{1}
\setcounter{secnumdepth}{1}
\setlist[itemize]{leftmargin=3em,itemsep=5pt,topsep=4pt,parsep=0pt}
\allowdisplaybreaks
\newcommand{\N}{\mathbb{N}}
\newcommand{\Z}{\mathbb{Z}}
\newcommand{\Q}{\mathbb{Q}}
\newcommand{\R}{\mathbb{R}}
\newcommand{\C}{\mathbb{C}}
\newcommand{\F}{\mathbb{F}}
\newcommand{\A}{\mathbb{A}}
\newcommand{\PP}{\mathbb P}
\newcommand{\E}{\mathbb{E}}
\newcommand{\eps}{\varepsilon}
\newcommand{\dd}{\,\mathrm d}
\newcommand{\sref}[2]{\hyperlink{source:#1:#2}{[S#2]}}
\newcommand{\eref}[2]{\hyperlink{extra:#1:#2}{[E#2]}}
% Turn the next switch off for a shorter reading copy. The quotations preserve
% every exactTarget field, including qualifications omitted from a short summary.
\newif\ifshowcatalogue
\showcataloguetrue
\newcommand{\cataloguescope}[1]{\ifshowcatalogue
 \paragraph{Exact scope in the supplied catalogue.}
 \begin{quote}\small #1\end{quote}\fi}
\usepackage{etoolbox}
\AtBeginEnvironment{itemize}{\small}
\numberwithin{equation}{section}
% Standalone export. Original section 118 preserved verbatim outside marked additions.
\begin{document}
\setcounter{section}{117}
% ================================================================
% problemId: problem.grothendieckkatz-p-curvature-conjecture
% source releaseRank: 118; source releaseStatus: open_with_solved_subcases
% exactTarget SHA-256: cc7704f0ad6a9b261ef34c1cbc1bb3eb390565067ba938e064e66d5a3d0871d0
\clearpage
\section{\texorpdfstring{Grothendieck–Katz p-Curvature Conjecture}{Grothendieck–Katz p-Curvature Conjecture}}\label{118}
\subsection{Definitions and mathematical statement}
For a connection in characteristic $p$, its $p$-curvature is the map on derivations
\[
 \psi_p(D)=\nabla_D^{\,p}-\nabla_{D^{[p]}},
\]
where $D^{[p]}$ is the $p$-fold derivation. Start with an algebraic integrable connection in characteristic zero and a model over a finitely generated $\Z$-algebra. If the $p$-curvature vanishes for almost all reductions in the source's sense, the conjecture predicts finite monodromy. Equivalently, after a finite étale cover the bundle with connection has a full horizontal basis. Vanishing at finitely many primes does not suffice, and a choice of model is part of expressing the arithmetic hypothesis.
\par\smallskip\textit{Formulation and concept references:} \sref{118}{1}.
\cataloguescope{Let X be a smooth connected complex variety and (V,\ensuremath{\nabla}) an algebraic vector bundle with integrable connection, spread out over a finitely generated Z-algebra. Prove that if the reductions have zero p-curvature for almost all primes, then (V,\ensuremath{\nabla}) has a full set of algebraic solutions, equivalently becomes trivial on a finite étale cover and has finite monodromy.}
\subsection{Short English statement}
Does vanishing $p$-curvature at almost every prime force a characteristic-zero differential equation to have only algebraic solutions and finite monodromy?
% BEGIN RESEARCH ADDITION 118
\subsection{Research attempt}

\paragraph{Current frontier and arithmetic scope.}
The source's Conjecture 1 quantifies over almost all maximal ideals of the
chosen finitely generated $\Z$-algebra, not one conveniently chosen
specialization in each residue characteristic. Its introduction records
Katz's regular-singularity theorem and the solvable-monodromy cases;
Theorem 3 proves the conjecture on specified higher-rank locally symmetric
varieties, with both rational-rank and real-rank hypotheses. Those hypotheses
do not include arbitrary curves. \sref{118}{1}

F\"urnsinn--Pannier's 2025 manuscript gives an effective, height- and
degree-dependent prime test for order-one equations. Lam--Litt's 2026
manuscript studies the isomonodromy foliation for smooth projective families
and its action on integral characters. These are relevant specialized
results, not a theorem for every integrable connection in this entry;
neither manuscript is used as an unverified black box below.
\eref{118}{1}\eref{118}{2}

\paragraph{Editorial clarification: where regular singularity enters.}
The original equivalence with finite analytic monodromy is valid here after
using the regular-singularity consequence of the arithmetic hypothesis.
Without that step it would be false: $d-dx$ on $\A^1$ has solution $e^x$ and
trivial topological monodromy, but no nonzero algebraic horizontal section.
For $\partial=d/dx$ in characteristic $p$, $\partial^{[p]}=0$ and
$(\partial-1)^p=-1$, so this example has nonzero $p$-curvature and is not a
counterexample to the stated conjecture. The role of regular singularity is
explicit in the source's introduction. \sref{118}{1}

\paragraph{Attack route.}
Try to convert the arithmetic vanishing into residue constraints, then
integrate those constraints to obtain an algebraic horizontal basis. On a
punctured projective line with a fixed nilpotent matrix this can be done
completely and explicitly. Two tests then probe the extension to the target:
finite local monodromy need not generate a finite group, and arbitrarily
many successful prime checks do not replace the almost-all-primes
hypothesis. These isolate the genuinely global obstruction rather than
assuming it disappears after the local calculation.

\paragraph{Attempt: a full calculation in a restricted unipotent class.}
Let $K$ be a number field,
$X=\PP^1_K\setminus S$, and let $\omega$ be a rational differential regular
on $X$. On the trivial rank-$r$ bundle consider
\begin{equation}\label{118:unipotent}
 \nabla=d-N\omega,\qquad
 N\in M_r(K),\quad N^r=0.
\end{equation}
The matrix is constant. Connections with noncommuting coefficient matrices
are not included. All models below exclude the finitely many bad primes
needed to define the data.

\textbf{Lemma.}
If \eqref{118:unipotent} has zero $p$-curvature at all but finitely many prime
ideals of $K$, then it is algebraically trivial on $X$. If $N\ne0$, this is
equivalent to $\omega=df$ for a rational function $f\in K(x)$ regular on $X$.

\emph{Proof.}
For $N=0$ there is nothing to prove. Suppose $N\ne0$. After a finite extension
$L/K$ splitting the poles, partial fractions and characteristic-zero
integration give
\begin{equation}\label{118:partial}
 \omega=df+\sum_{i=1}^s r_i\frac{dx}{x-\alpha_i},
 \qquad f\in L(x),\quad r_i\in L.
\end{equation}
The residue at infinity is minus the sum of the displayed residues. Polynomial
terms and every pole term of order greater than one have a rational primitive;
this proves \eqref{118:partial}, including the possible pole at infinity.
The arithmetic hypothesis persists under the finite extension.

The finite exponential
\begin{equation}\label{118:gauge}
 F=\exp(Nf)=\sum_{j=0}^{r-1}\frac{N^jf^j}{j!},
 \qquad F^{-1}=\exp(-Nf),
\end{equation}
has determinant one. Since $N$ commutes with $F$ and $dF=N\,df\,F$,
\[
 F^{-1}\nabla F=d-N(\omega-df).
\]
This is a rational gauge transformation. It is defined on $X$ and in every
sufficiently good reduction, after clearing the finite set of denominators
arising from $f,N$ and the factorials. Vanishing of $p$-curvature is
preserved by gauge transformation.

Near the pole $\alpha_i$, put $z=x-\alpha_i$ and
$\theta=z\,d/dz$. The transformed logarithmic connection satisfies
\[
 \nabla_\theta=\theta-N(r_i+zh(z)),\qquad h(z)\in L[\![z]\!].
\]
This operator preserves the logarithmic lattice and acts on its fiber
at $z=0$ as $-r_iN$. In characteristic $p$, $\theta^{[p]}=\theta$;
this follows directly from $\theta(z^m)=mz^m$ and $m^p=m$ in the residue
field. Therefore the fiber of its $p$-curvature is
\begin{equation}\label{118:residue}
 \psi_p(\theta)\big|_{z=0}=(-r_iN)^p-(-r_iN)=r_iN,
 \qquad p\ge r.
\end{equation}
Although the point $z=0$ is removed from $X$, the logarithmic operator extends
there. If its $p$-curvature vanishes on the punctured neighborhood, it vanishes
as a rational matrix and hence also on this fiber. Thus \eqref{118:residue}
is zero at almost all prime ideals of $L$.

A nonzero entry of $N$ stays nonzero modulo all but finitely many prime
ideals. A nonzero algebraic number $r_i$, after denominators are cleared,
also vanishes modulo only finitely many prime ideals: any such prime divides
its nonzero integral norm. It follows that every $r_i=0$ in characteristic
zero. Hence $\omega=df$. Taking the trace of $f$ from $L(x)$ to $K(x)$ and
dividing by $[L:K]$ gives a primitive in $K(x)$ itself.

A rational primitive has no pole on $X$, because in characteristic zero a
pole of order $m>0$ of $f$ gives a pole of order $m+1$ of $df$. Thus
\eqref{118:gauge} and its inverse are regular on $X$. Their columns form a
full horizontal basis. Conversely this rational gauge to the trivial
connection gives zero $p$-curvature in every good reduction. $\square$

The lemma treats arbitrary fixed nilpotence index, not just a $2\times2$
example, but it is a special solvable class already compatible with the
known literature. No novelty claim is made. The point of the proof is to
expose every transition from reduction modulo primes to a characteristic-zero
primitive, without inferring a global theorem from finitely many computations.
\sref{118}{1}

\paragraph{First obstruction: residue tests lose global information.}
The argument used that a rational differential on $\PP^1$ with zero residues
has a rational primitive. That assertion fails on a curve of positive genus.
For example, on an elliptic curve a nonzero regular differential has no
residues, but cannot equal $df$ for a rational function: such an $f$ could
have no poles, hence would be constant on the projective curve. Thus one
cannot replace the projective line in the lemma by an arbitrary curve using
only its residue computation. This does \emph{not} assert that the resulting
connection satisfies the almost-all-primes condition.

Nor does finite order of all boundary generators imply finite monodromy.
The rational matrices
\begin{equation}\label{118:matrices}
 A=\begin{pmatrix}0&-1\\1&0\end{pmatrix},\qquad
 B=\begin{pmatrix}0&-4\\1/4&0\end{pmatrix}
\end{equation}
satisfy $A^2=B^2=-I$, whereas
\[
 AB=\operatorname{diag}(-1/4,-4)
\]
has infinite order. On the abstract fundamental group of a four-punctured
sphere, assigning its boundary generators to $A,B,B^{-1},A^{-1}$ respects
the product-one relation. Each boundary image is finite order while the
image group is infinite. This is a representation-theoretic obstruction to
a proposed inference, not a connection claimed to have vanishing
$p$-curvatures.

The source's locally symmetric theorem gets around this obstruction through
normal-subgroup rigidity under its rank hypotheses. There is no corresponding
argument here for general fundamental groups. Removing those hypotheses
from the cited theorem would assume the difficult missing step.
\sref{118}{1}

\paragraph{Second obstruction: arbitrarily many good tests can be misleading.}
For a finite set $\mathcal P$ of rational primes, let
$c=\prod_{\ell\in\mathcal P}\ell$ and on $\mathbb G_m$ take
\[
 N=\begin{pmatrix}0&1\\0&0\end{pmatrix},\qquad
 \nabla=d-cN\frac{dx}{x}.
\]
Since $\theta=x\partial_x$ commutes with $N$, $N^p=0$ for every $p\ge2$, and
$\theta^{[p]}=\theta$, the exact calculation is
\begin{equation}\label{118:finite-test}
 \psi_p(\theta)=(\theta-cN)^p-(\theta-cN)=cN.
\end{equation}
It vanishes precisely at the selected primes. Frobenius-semilinearity of
$p$-curvature shows that vanishing for the generator $\theta$ is vanishing
for every derivation on this curve. In characteristic zero a horizontal
fundamental matrix is $I+cN\log x$, and a loop around zero has monodromy
\[
 I+2\pi i cN,
 \qquad (I+2\pi i cN)^k=I+2\pi i k cN.
\]
The monodromy is infinite. The example defeats any height-independent fixed
finite-prime verification rule for all these rank-two systems, not the
almost-all-primes hypothesis. Its coefficient height grows with the prime
set; it does not conflict with the height-dependent order-one algorithm
reported in the 2025 manuscript. \eref{118}{1}

\paragraph{Stress test.}
The source's arithmetic quantifier is preserved throughout the discussion
of the target. The proved lemma intentionally restricts to a number field,
a punctured projective line, a trivial bundle, and a single constant nilpotent
matrix multiplying one differential. It does not settle arbitrary
finitely generated models, higher-genus curves, nontrivial bundles, or
noncommuting monodromy. The gauges are algebraic and invertible on the
stated open set; a merely formal or analytic horizontal solution is not
being passed off as an algebraic one.

The finite-order matrices, the elliptic differential, and the finite-prime
family each disprove a separate proposed shortcut. None supplies a
counterexample satisfying all the arithmetic hypotheses of the conjecture.
The companion script checks the nilpotent gauge identity, the finite-order
matrix products, and the selected-prime calculation exactly.

\paragraph{Outcome.}
\textbf{Outcome: Exploratory approach with unresolved gap.}

A full, explicit proof is supplied for the fixed-nilpotent, one-differential
class \eqref{118:unipotent} on punctured projective lines, together with a
family passing any prescribed finite prime test but having infinite
monodromy. These are a restricted theorem and an obstruction to a verification
shortcut, respectively, not a solution or counterexample to the general
Grothendieck--Katz conjecture.

The remaining problem is global: obtain finite monodromy from the arithmetic
hypothesis without special solvable structure or higher-rank rigidity.
The residue calculation cannot see enough of that global information, and
neither local finite order nor a finite collection of reductions supplies it.
% END RESEARCH ADDITION 118
\subsection{Sources}
\begin{itemize}\raggedright
\item[\textnormal{[S1]}]\hypertarget{source:118:1}{} Benson Farb and Mark Kisin, The Grothendieck–Katz Conjecture for Certain Locally Symmetric Varieties.
\url{https://www.math.uchicago.edu/~farb/papers/pcurv.pdf}.
{\small\textit{Catalogue formulation source.}}
% BEGIN NEW REFERENCES 118
\item[\textnormal{[E1]}]\hypertarget{extra:118:1}{} Florian F\"urnsinn and Lucas Pannier, \emph{An Effective Version of the $p$-Curvature Conjecture for Order One Differential Equations}, arXiv:2510.00892v1, 1 October 2025, abstract and introduction. Height- and degree-dependent effective prime tests; manuscript result, not used as a black box.
\url{https://arxiv.org/html/2510.00892v1}.
\item[\textnormal{[E2]}]\hypertarget{extra:118:2}{} Yeuk Hay Joshua Lam and Daniel Litt, \emph{$p$-curvature and non-abelian cohomology}, arXiv:2601.07933v1, 12 January 2026, abstract and Section 1.1. Isomonodromy foliations and integral characters; not the unrestricted target of this entry.
\url{https://arxiv.org/html/2601.07933v1}.
% END NEW REFERENCES 118
\end{itemize}

\end{document}
